%! Change in Arithmetic and Geometric Sequences — Unit 4, Lesson 4.1 — Guided Notes (Answer Key)
\documentclass[10pt]{article}
\usepackage{precalculus-article}
\usepackage{precalculus-key}   % pulls in -boxes; do NOT also load -boxes
\newcommand{\TallMath}[1]{$\displaystyle #1\rule[-1.4em]{0pt}{3.2em}$}

\begin{document}

\pageheader{Unit 4, Lesson 4.1}{Guided Notes --- Answer Key}
\namedateperiod

\vspace{0.06in}

\begin{objectivebox}
By the end of this lesson, I will be able to\ldots
\begin{enumerate}[itemsep=2pt]
  \item Express an \ans{arithmetic} sequence as a function $a_n = a_0 + dn$ and identify the
        common \ans{difference}. \hfill\textit{(LO 2.1.A)}
  \item Express a \ans{geometric} sequence as a function $g_n = g_0 r^n$ and identify the
        common \ans{ratio}. \hfill\textit{(LO 2.1.B)}
  \item Compare how arithmetic and geometric sequences \ans{change} across successive terms. \hfill\textit{(LO 2.1.B)}
\end{enumerate}
\end{objectivebox}

\vspace{0.04in}

\begin{vocabbox}
\termblanklong{sequence}
\termblanklong{arithmetic sequence}
\termblanklong{geometric sequence}
\termblanklong{common difference}
\termblanklong{common ratio}
\termblanklong{general term (explicit formula)}
\termblanklong{recursive formula}
\end{vocabbox}

\vspace{0.04in}

\begin{hookbox}
Two sequences are written on the board:
\[
  2,\ 5,\ 8,\ 11,\ \ldots \qquad\qquad 2,\ 6,\ 18,\ 54,\ \ldots
\]
The first sequence \ans{increases} by the same amount each step.
The second sequence \ans{increases} by the same factor each step.

\medskip
Estimate: how many steps until each sequence first exceeds 1000?

\medskip
Sequence 1: about \ans{333} steps. \quad Sequence 2: about \ans{6} steps.
\end{hookbox}

\vspace{0.04in}

\begin{notesbox}{LO 2.1.A --- Arithmetic Sequences}
\textbf{Definition:} A \textit{sequence} is a function from the \ans{whole} numbers to the real numbers.
Its graph consists of \ans{discrete} points (not a curve).

\medskip
\textbf{Arithmetic sequence:} Successive terms differ by a \ans{constant} value $d$, the
\textit{common difference}.

\[
  d = a_n - a_{n-1} \quad \text{(constant for all } n \geq 1\text{)}
\]

\textbf{General (explicit) formula:}
\[
  a_n = a_0 + d \cdot n
\]
where $a_0$ is the \ans{initial} value (the term when $n = 0$).

\medskip
\textbf{Alternate form:} If we know the $k$th term $a_k$:
\[
  a_n = a_k + d(n - k)
\]

\medskip
\textbf{Example:} The sequence $2,\ 5,\ 8,\ 11,\ \ldots$

\smallskip
Common difference: $d = $ \ans{3}

\smallskip
Explicit formula: $a_n = $ \ans{$2 + 3n$}

\smallskip
Find the first term exceeding 1000: \ans{$2 + 3n > 1000 \Rightarrow n > 332.6$, so $n = 333$}
\end{notesbox}

\vspace{0.04in}

\begin{notesbox}{LO 2.1.B --- Geometric Sequences}
\textbf{Geometric sequence:} Successive terms have a \ans{constant} ratio $r$, the
\textit{common ratio}.

\[
  r = \frac{g_n}{g_{n-1}} \quad \text{(constant for all } n \geq 1\text{)}
\]

\textbf{General (explicit) formula:}
\[
  g_n = g_0 \cdot r^n
\]
where $g_0$ is the \ans{initial} value.

\medskip
\textbf{Alternate form:} If we know the $k$th term $g_k$:
\[
  g_n = g_k \cdot r^{(n-k)}
\]

\medskip
\textbf{Example:} The sequence $2,\ 6,\ 18,\ 54,\ \ldots$

\smallskip
Common ratio: $r = $ \ans{3}

\smallskip
Explicit formula: $g_n = $ \ans{$2 \cdot 3^n$}

\smallskip
Find the first term exceeding 1000: \ans{$2 \cdot 3^n > 1000 \Rightarrow 3^n > 500$; $3^6 = 729 > 500$, so $n = 6$}
\end{notesbox}

\vspace{0.04in}

\begin{notesbox}{Comparison: Arithmetic vs.\ Geometric Change (EK 2.1.B.3)}
Complete the table for both sequences starting at $n = 0$:

\renewcommand{\arraystretch}{1.8}
\begin{tabular}{|c|c|c|c|}
\hline
$n$ & Arithmetic: $a_n = 2 + 3n$ & Geometric: $g_n = 2 \cdot 3^n$ & Difference $g_n - a_n$ \\
\hline
0 & 2 & 2 & 0 \\
\hline
1 & \ans{5}  & \ans{6}   & \ans{1}    \\
\hline
2 & \ans{8}  & \ans{18}  & \ans{10}   \\
\hline
3 & \ans{11} & \ans{54}  & \ans{43}   \\
\hline
4 & \ans{14} & \ans{162} & \ans{148}  \\
\hline
5 & \ans{17} & \ans{486} & \ans{469}  \\
\hline
\end{tabular}

\medskip
\textbf{Key insight:} Arithmetic sequences add \ans{a constant} each step;
geometric sequences multiply by \ans{a constant} each step.
Increasing geometric sequences eventually grow \ans{faster} than arithmetic sequences.
\end{notesbox}

\end{document}
